\documentclass[12pt]{beamer}
\usetheme{Madrid}
\usecolortheme{default}
\usepackage{amsmath}
\usepackage{amssymb}

\title{Solving Other Equations}
\subtitle{Absolute Value and Rational Equations}
\date{}

\begin{document}
\setbeamertemplate{navigation symbols}{}
\setbeamertemplate{footline}[frame number]

\frame{\titlepage}

% ==================== DO NOW ====================
\begin{frame}{Do Now}
\textbf{Solve the following problems:}

\vspace{0.5cm}

\textbf{1.} Solve: $x^2 = 49$

\vspace{0.5cm}

\textbf{2.} Solve: $(x-3)^2 = 16$

\vspace{0.5cm}

\textbf{3.} Solve: $x^2 + 4x - 5 = 0$

\vspace{0.5cm}

\textbf{4.} A ball is dropped from a height of 64 feet. When does it hit the ground?
$$h(t) = -16t^2 + v_0t + h_0$$
\end{frame}

\begin{frame}{Do Now - Answers}
\textbf{1.} $x^2 = 49$ \quad $\Rightarrow$ \quad $x = 7$ or $x = -7$

\vspace{0.5cm}

\textbf{2.} $(x-3)^2 = 16$ \quad $\Rightarrow$ \quad $x = 7$ or $x = -1$

\vspace{0.5cm}

\textbf{3.} $x^2 + 4x - 5 = 0$ \quad $\Rightarrow$ \quad $x = 1$ or $x = -5$

\vspace{0.5cm}

\textbf{4.} $h(t) = -16t^2 + 0 \cdot t + 64 = 0$

$-16t^2 + 64 = 0 \quad \Rightarrow \quad t^2 = 4 \quad \Rightarrow \quad t = 2$ seconds
\end{frame}

% ==================== PURPOSE ====================
\begin{frame}{Today's Goal}
\Large
\textbf{By the end of today, you will be able to:}

\vspace{0.5cm}

\normalsize
\begin{itemize}
\item Solve absolute value equations
\item Solve rational equations
\end{itemize}

\vspace{0.5cm}

\textbf{Why?} These techniques build on linear equation solving and prepare you for more complex algebraic problem-solving.
\end{frame}

% ==================== DIAGNOSTICS - LINEAR EQUATIONS ====================
\begin{frame}{Diagnostic: Linear Equations}
\textbf{Solve:} $3(2x - 5) = 4x + 7$

\vspace{1cm}

\textbf{A.} $x = 11$ \hfill \uncover<2->{{\color{green}\checkmark}}

\textbf{B.} $x = -4$

\textbf{C.} $x = \displaystyle\frac{22}{2}$

\textbf{D.} No solution
\end{frame}

\begin{frame}{Diagnostic: Linear Equations (Follow-up)}
\textbf{Solve:} $2(3x + 4) = 5x - 1$

\vspace{1cm}

\textbf{A.} $x = -9$ \hfill \uncover<2->{{\color{green}\checkmark}}

\textbf{B.} $x = 7$

\textbf{C.} $x = \displaystyle\frac{-9}{1}$

\textbf{D.} No solution
\end{frame}

\begin{frame}{Diagnostic: Linear Equations}
\textbf{Solve:} $4(x + 2) = 4x - 5$

\vspace{1cm}

\textbf{A.} $x = 0$

\textbf{B.} $x = -13$

\textbf{C.} $x = \displaystyle\frac{-13}{0}$

\textbf{D.} No solution \hfill \uncover<2->{{\color{green}\checkmark}}
\end{frame}

\begin{frame}{Diagnostic: Linear Equations (Follow-up)}
\textbf{Solve:} $5(x - 3) = 5x + 7$

\vspace{1cm}

\textbf{A.} $x = 0$

\textbf{B.} $x = 22$

\textbf{C.} $x = \displaystyle\frac{22}{0}$

\textbf{D.} No solution \hfill \uncover<2->{{\color{green}\checkmark}}
\end{frame}

% ==================== SKILL 1: ABSOLUTE VALUE - CYCLE 1 ====================
\begin{frame}{Skill 1: Absolute Value Equations}
\Large
\textbf{Cycle 1: Isolate and Solve}
\end{frame}

\begin{frame}{Atom 1 Diagnostic: Isolating}
\textbf{If} $|x + 5| - 3 = 7$, \textbf{what does} $|x + 5|$ \textbf{equal?}

\vspace{1cm}

\textbf{A.} 4

\textbf{B.} 7

\textbf{C.} 10 \hfill \uncover<2->{{\color{green}\checkmark}}

\textbf{D.} $-3$
\end{frame}

\begin{frame}{Atom 1 Teaching: Isolating the Absolute Value}
$|x + 2| - 3 = 6$ \pause $\rightarrow$ $|x + 2| = 9$

\pause
\vspace{0.5cm}

$|x + 2| + 5 = 12$ \pause $\rightarrow$ $|x + 2| = 7$

\pause
\vspace{0.5cm}

$3|x + 2| = 6$ \pause $\rightarrow$ $|x + 2| = 2$

\pause
\vspace{0.5cm}

$\displaystyle\frac{|x + 2|}{4} = 3$ \pause $\rightarrow$ $|x + 2| = 12$
\end{frame}

\begin{frame}{Atom 2 Teaching: Absolute Value Property}
$|x| = 5$ \pause $\rightarrow$ $x = 5$ or $x = -5$

\pause
\vspace{0.5cm}

$|x - 1| = 8$ \pause $\rightarrow$ $x - 1 = 8$ or $x - 1 = -8$

\pause
\vspace{0.5cm}

$|x + 2| = 7$ \pause $\rightarrow$ $x + 2 = 7$ or $x + 2 = -7$
\end{frame}

\begin{frame}{Atom 2 Test 1}
\textbf{If} $|y| = 9$, \textbf{what are the valid solutions?}

\vspace{1cm}

\textbf{A.} $y = 9$ only

\textbf{B.} $y = 9$ or $y = -9$ \hfill \uncover<2->{{\color{green}\checkmark}}

\textbf{C.} $y = 9$ and $y = 9$

\textbf{D.} $y - 9 = 0$ or $y + 9 = 0$
\end{frame}

\begin{frame}{Atom 2 Test 2}
\textbf{If} $|x - 3| = 6$, \textbf{which equations should we solve?}

\vspace{1cm}

\textbf{A.} $x = 6$ or $x = -6$

\textbf{B.} $x - 3 = 6$ only

\textbf{C.} $x - 3 = 6$ or $x - 3 = -6$ \hfill \uncover<2->{{\color{green}\checkmark}}

\textbf{D.} $x = 3$ or $x = -3$
\end{frame}

\begin{frame}{I Do: Absolute Value}
\textbf{Solve:} $|2x - 5| + 3 = 10$

\pause
\vspace{0.3cm}

\textbf{Step 1: Isolate} $\Rightarrow$ $|2x - 5| = 7$

\pause
\vspace{0.3cm}

\textbf{Step 2: Two cases} $\Rightarrow$ $2x - 5 = 7$ or $2x - 5 = -7$

\pause
\vspace{0.3cm}

\textbf{Step 3: Solve each case}

\pause
Case 1: $2x - 5 = 7 \Rightarrow 2x = 12 \Rightarrow x = 6$

\pause
Case 2: $2x - 5 = -7 \Rightarrow 2x = -2 \Rightarrow x = -1$

\pause
\vspace{0.3cm}

\textbf{Step 4: Check both solutions}

\pause
Check $x = 6$: $|2(6) - 5| + 3 = |7| + 3 = 10$ \checkmark

\pause
Check $x = -1$: $|2(-1) - 5| + 3 = |-7| + 3 = 10$ \checkmark

\pause
\vspace{0.3cm}

\Large \textbf{Answer:} $x = 6$ or $x = -1$
\end{frame}

\begin{frame}{We Do: Absolute Value}
\textbf{Solve:} $|3x + 1| - 2 = 5$

\vspace{2cm}

\textit{Work through this problem step by step.}
\end{frame}

\begin{frame}{We Do: Absolute Value - Solution}
\textbf{Solve:} $|3x + 1| - 2 = 5$

\pause
\vspace{0.3cm}

\textbf{Step 1:} Isolate \pause $\Rightarrow$ $|3x + 1| = 7$

\pause
\vspace{0.3cm}

\textbf{Step 2:} Two cases \pause $\Rightarrow$ $3x + 1 = 7$ or $3x + 1 = -7$

\pause
\vspace{0.3cm}

\textbf{Step 3:} Solve:

\pause
\begin{itemize}
\item Case 1: $3x = 6 \Rightarrow x = 2$ \pause
\item Case 2: $3x = -8 \Rightarrow x = -\displaystyle\frac{8}{3}$
\end{itemize}

\pause
\vspace{0.3cm}

\textbf{Step 4:} Check both (both work!) \checkmark

\pause
\vspace{0.3cm}

\textbf{Answer:} $x = 2$ or $x = -\displaystyle\frac{8}{3}$
\end{frame}

\begin{frame}{Practice: Absolute Value}
\textbf{Solve the following:}

\vspace{0.5cm}

\textbf{1.} $|x - 4| + 6 = 11$ \uncover<2->{\hfill $x = 9$ or $x = -1$}

\vspace{0.5cm}

\textbf{2.} $|2x + 3| - 1 = 8$ \uncover<3->{\hfill $x = 3$ or $x = -6$}

\vspace{0.5cm}

\textbf{3.} $|5x - 2| + 4 = 12$ \uncover<4->{\hfill $x = 2$ or $x = -\displaystyle\frac{6}{5}$}
\end{frame}

% ==================== SKILL 1: ABSOLUTE VALUE - CYCLE 2 ====================
\begin{frame}{Skill 1: Absolute Value Equations}
\Large
\textbf{Cycle 2: No-Solution Case}
\end{frame}

\begin{frame}{Atom 4 Teaching: Absolute Value is Always Non-Negative}
$|x| = -5$ \pause \quad {\color{red}NO}

\pause
\vspace{0.5cm}

$|x| = 3$ \pause \quad {\color{green}YES}

\pause
\vspace{0.5cm}

$|x| = 0$ \pause \quad {\color{green}YES}

\pause
\vspace{0.5cm}

$|x + 2| = 8$ \pause \quad {\color{green}YES}

\pause
\vspace{0.5cm}

$|3x - 1| = -2$ \pause \quad {\color{red}NO}
\end{frame}

\begin{frame}{Atom 4 Test 1}
\textbf{Does} $|x - 5| = 7$ \textbf{have solutions?}

\vspace{1cm}

\textbf{A.} Yes \hfill \uncover<2->{{\color{green}\checkmark}}

\textbf{B.} No
\end{frame}

\begin{frame}{Atom 4 Test 2}
\textbf{Does} $|2x + 3| = -4$ \textbf{have solutions?}

\vspace{1cm}

\textbf{A.} Yes

\textbf{B.} No \hfill \uncover<2->{{\color{green}\checkmark}}
\end{frame}

\begin{frame}{Atom 4 Test 3}
\textbf{Does} $|x| = 0$ \textbf{have solutions?}

\vspace{1cm}

\textbf{A.} Yes \hfill \uncover<2->{{\color{green}\checkmark}}

\textbf{B.} No
\end{frame}

\begin{frame}{I Do: No-Solution Case}
\textbf{Solve:} $|3x + 1| + 7 = 4$

\pause
\vspace{0.5cm}

\textbf{Step 1: Isolate} $\Rightarrow$ $|3x + 1| = -3$

\pause
\vspace{0.5cm}

\textbf{Step 2: Recognize the problem}

Absolute value cannot equal a negative number!

\pause
\vspace{0.5cm}

\Large
\textbf{Answer: No solution}
\end{frame}

\begin{frame}{We Do: No-Solution Case}
\textbf{Solve:} $|2x - 5| + 9 = 6$

\vspace{2cm}

\textit{Work through this problem.}
\end{frame}

\begin{frame}{We Do: No-Solution Case - Solution}
\textbf{Solve:} $|2x - 5| + 9 = 6$

\pause
\vspace{0.5cm}

\textbf{Step 1:} Isolate \pause $\Rightarrow$ $|2x - 5| = -3$

\pause
\vspace{0.5cm}

\textbf{Step 2:} Absolute value cannot be negative!

\pause
\vspace{0.5cm}

\Large
\textbf{Answer: No solution}
\end{frame}

% ==================== SKILL 2: RATIONAL EQUATIONS ====================
\begin{frame}{Skill 2: Rational Equations}
\Large
\textbf{Solving with the LCD Method}
\end{frame}

\begin{frame}{Atom 1 Diagnostic: Finding LCD}
\textbf{What is the LCD of} $\displaystyle\frac{1}{4}$, $\displaystyle\frac{1}{6}$, \textbf{and} $\displaystyle\frac{1}{8}$?

\vspace{1cm}

\textbf{A.} 12

\textbf{B.} 24 \hfill \uncover<2->{{\color{green}\checkmark}}

\textbf{C.} 48

\textbf{D.} 192
\end{frame}

\begin{frame}{Atom 1 Diagnostic Follow-up}
\textbf{What is the LCD of} $\displaystyle\frac{1}{3}$, $\displaystyle\frac{1}{5}$, \textbf{and} $\displaystyle\frac{1}{15}$?

\vspace{1cm}

\textbf{A.} 15 \hfill \uncover<2->{{\color{green}\checkmark}}

\textbf{B.} 45

\textbf{C.} 75

\textbf{D.} 225
\end{frame}

\begin{frame}{Atom 2 Diagnostic: LCD with Variables}
\textbf{What is the LCD of} $\displaystyle\frac{1}{2x}$, $\displaystyle\frac{1}{3x}$, \textbf{and} $\displaystyle\frac{1}{6}$?

\vspace{1cm}

\textbf{A.} $6x$ \hfill \uncover<2->{{\color{green}\checkmark}}

\textbf{B.} $6x^2$

\textbf{C.} $36x$

\textbf{D.} $2x$
\end{frame}

\begin{frame}{Atom 2 Diagnostic Follow-up}
\textbf{What is the LCD of} $\displaystyle\frac{1}{4x}$, $\displaystyle\frac{1}{5x}$, \textbf{and} $\displaystyle\frac{1}{10}$?

\vspace{1cm}

\textbf{A.} $20x$ \hfill \uncover<2->{{\color{green}\checkmark}}

\textbf{B.} $10x$

\textbf{C.} $20x^2$

\textbf{D.} $40x$
\end{frame}

\begin{frame}{Atom 3 Diagnostic: Multiply by LCD}
\textbf{If we multiply both sides of} $\displaystyle\frac{1}{3} + \frac{1}{x} = \frac{1}{2}$ \textbf{by} $6x$, \textbf{what is the simplified equation?}

\vspace{0.5cm}

\textbf{A.} $2x + 6 = 3x$ \hfill \uncover<2->{{\color{green}\checkmark}}

\textbf{B.} $2 + 6x = 3$

\textbf{C.} $\displaystyle\frac{6x}{3} + \frac{6x}{x} = \frac{6x}{2}$

\textbf{D.} $3x + 2 = 6x$
\end{frame}

\begin{frame}{Atom 3 Diagnostic Follow-up}
\textbf{If we multiply both sides of} $\displaystyle\frac{1}{4} - \frac{1}{x} = \frac{1}{3}$ \textbf{by} $12x$, \textbf{what is the simplified equation?}

\vspace{0.5cm}

\textbf{A.} $3x - 12 = 4x$ \hfill \uncover<2->{{\color{green}\checkmark}}

\textbf{B.} $3 - 12x = 4$

\textbf{C.} $12 - 3x = 4x$

\textbf{D.} $\displaystyle\frac{12x}{4} - \frac{12x}{x} = \frac{12x}{3}$
\end{frame}

\begin{frame}{Atom 6 Test 1: Domain Restrictions}
\textbf{For the equation} $\displaystyle\frac{1}{x-2} + \frac{1}{x+5} = 3$, \textbf{which values of} $x$ \textbf{are NOT allowed?}

\vspace{0.5cm}

\textbf{A.} $x = 2$ only

\textbf{B.} $x = -5$ only

\textbf{C.} $x = 2$ or $x = -5$ \hfill \uncover<2->{{\color{green}\checkmark}}

\textbf{D.} $x = 0$
\end{frame}

\begin{frame}{Atom 6 Test 2: Checking Solutions}
\textbf{You solved a rational equation and got} $x = 3$. \textbf{The original equation was} $\displaystyle\frac{2}{x-3} + \frac{1}{x} = 5$. \textbf{Is} $x = 3$ \textbf{a valid solution?}

\vspace{0.5cm}

\textbf{A.} Yes

\textbf{B.} No, because it makes a denominator zero \hfill \uncover<2->{{\color{green}\checkmark}}

\textbf{C.} Yes, because $3 \neq 0$

\textbf{D.} Cannot determine
\end{frame}

\begin{frame}{I Do 1: Rational Equations}
\textbf{Solve:} $\displaystyle\frac{1}{2} + \frac{1}{3} = \frac{1}{x}$

\pause
\vspace{0.3cm}

\textbf{Step 1:} LCD = $6x$ \quad \pause \textbf{Step 2:} Multiply by LCD

\pause
$6x \cdot \displaystyle\frac{1}{2} + 6x \cdot \frac{1}{3} = 6x \cdot \frac{1}{x}$

\pause
\vspace{0.3cm}

\textbf{Step 3:} Simplify $\Rightarrow$ $3x + 2x = 6$

\pause
\vspace{0.3cm}

\textbf{Step 4:} Solve $\Rightarrow$ $5x = 6$ \pause $\Rightarrow$ $x = \displaystyle\frac{6}{5}$

\pause
\vspace{0.3cm}

\textbf{Step 5:} Check: $\displaystyle\frac{1}{2} + \frac{1}{3} = \frac{3}{6} + \frac{2}{6} = \frac{5}{6}$ and $\frac{1}{6/5} = \frac{5}{6}$ \checkmark

\pause
\vspace{0.3cm}

\Large \textbf{Answer:} $x = \displaystyle\frac{6}{5}$
\end{frame}

\begin{frame}{I Do 2: Rational Equations}
\textbf{Solve:} $\displaystyle\frac{x}{x-1} - \frac{6}{x+3} = \frac{x^2}{x^2+2x-3}$

\pause
\vspace{0.3cm}

\textbf{Step 1:} Factor right side $\Rightarrow$ $x^2 + 2x - 3 = (x-1)(x+3)$

\pause
$$\frac{x}{x-1} - \frac{6}{x+3} = \frac{x^2}{(x-1)(x+3)}$$

\pause
\vspace{0.3cm}

\textbf{Step 2:} LCD = $(x-1)(x+3)$

\pause
\vspace{0.3cm}

\textbf{Step 3:} Multiply by LCD and simplify $\Rightarrow$ $x(x+3) - 6(x-1) = x^2$
\end{frame}

\begin{frame}{I Do 2: Rational Equations (continued)}
$$x(x+3) - 6(x-1) = x^2$$

\pause
\vspace{0.3cm}

\textbf{Step 4:} Expand $\Rightarrow$ $x^2 + 3x - 6x + 6 = x^2$

\pause
$\Rightarrow$ $x^2 - 3x + 6 = x^2$

\pause
\vspace{0.3cm}

\textbf{Step 5:} Simplify (notice $x^2$ cancels!) $\Rightarrow$ $-3x + 6 = 0$

\pause
$\Rightarrow$ $-3x = -6$ \pause $\Rightarrow$ $x = 2$

\pause
\vspace{0.3cm}

\textbf{Step 6:} Check: $\displaystyle\frac{2}{1} - \frac{6}{5} = \frac{10}{5} - \frac{6}{5} = \frac{4}{5}$ and $\frac{4}{5} = \frac{4}{5}$ \checkmark

\pause
\vspace{0.3cm}

\Large \textbf{Answer:} $x = 2$
\end{frame}

\begin{frame}{We Do: Rational Equations (Extraneous Solution)}
\textbf{Solve:} $\displaystyle\frac{3}{x} - \frac{2}{x-3} = \frac{-12}{x^2-9}$

\vspace{2cm}

\textit{Work through this problem step by step. Be careful - one solution might be extraneous!}
\end{frame}

\begin{frame}{We Do: Solution}
$\displaystyle\frac{3}{x} - \frac{2}{x-3} = \frac{-12}{x^2-9}$

\pause
\vspace{0.3cm}

\textbf{Step 1:} Factor $\Rightarrow$ $x^2 - 9 = (x-3)(x+3)$ \quad LCD = $x(x-3)(x+3)$

\pause
$$\frac{3}{x} - \frac{2}{x-3} = \frac{-12}{(x-3)(x+3)}$$

\pause
\vspace{0.3cm}

\textbf{Step 2:} Multiply by LCD, simplify $\Rightarrow$ $3(x-3)(x+3) - 2x(x+3) = -12x$

\pause
\vspace{0.3cm}

\textbf{Step 3:} Expand $\Rightarrow$ $3x^2 - 27 - 2x^2 - 6x = -12x$

\pause
$\Rightarrow$ $x^2 + 6x - 27 = 0$

\pause
\vspace{0.3cm}

\textbf{Step 4:} Factor $\Rightarrow$ $(x + 9)(x - 3) = 0$ \pause $\Rightarrow$ $x = -9$ or $x = 3$
\end{frame}

\begin{frame}{We Do: Solution (continued)}
$x = -9$ or $x = 3$

\pause
\vspace{0.3cm}

\textbf{Step 5: Check BOTH solutions}

\pause
\vspace{0.3cm}

\textbf{Check} $x = 3$: $\displaystyle\frac{3}{3} - \frac{2}{3-3} = \frac{-12}{9-9}$ \pause $\Rightarrow$ $\frac{3}{3} - \frac{2}{0} = \frac{-12}{0}$

\pause
\textbf{Division by zero!} This solution is {\color{red}EXTRANEOUS}.

\pause
\vspace{0.3cm}

\textbf{Check} $x = -9$: $\displaystyle\frac{3}{-9} - \frac{2}{-12} = -\frac{1}{3} + \frac{1}{6}$

\pause
$= -\displaystyle\frac{2}{6} + \frac{1}{6} = -\frac{1}{6}$

\pause
And $\displaystyle\frac{-12}{72} = -\frac{1}{6}$ \checkmark

\pause
\vspace{0.3cm}

\Large \textbf{Answer:} $x = -9$ only
\end{frame}

% ==================== CONSOLIDATION ====================
\begin{frame}{Consolidation: Practice Problems}
\textbf{Solve the following rational equations:}

\vspace{0.5cm}

\textbf{1.} $\displaystyle\frac{2}{x} + \frac{3}{x+1} = 5$

\vspace{0.5cm}

\textbf{2.} $\displaystyle\frac{x}{x+2} - \frac{4}{x-2} = \frac{16}{x^2-4}$

\vspace{0.5cm}

\textbf{3.} $\displaystyle\frac{5}{x} - \frac{3}{x-2} = \frac{-10}{x^2-2x}$
\end{frame}

\begin{frame}{Consolidation: Answers}
\textbf{1.} $\displaystyle\frac{2}{x} + \frac{3}{x+1} = 5$

\quad $\Rightarrow$ \quad $x = 1$ or $x = -\displaystyle\frac{1}{5}$

\vspace{0.5cm}

\textbf{2.} $\displaystyle\frac{x}{x+2} - \frac{4}{x-2} = \frac{16}{x^2-4}$

\quad $\Rightarrow$ \quad $x = 4$

\vspace{0.5cm}

\textbf{3.} $\displaystyle\frac{5}{x} - \frac{3}{x-2} = \frac{-10}{x^2-2x}$

\quad $\Rightarrow$ \quad $x = 5$ (Note: $x = 0$ is extraneous)
\end{frame}

% ==================== CLOSURE ====================
\begin{frame}{Closure}
\Large
\textbf{What We Learned Today:}

\vspace{0.5cm}

\normalsize
\textbf{Absolute Value Equations:}
\begin{itemize}
\item Isolate the absolute value first
\item Set up two cases: $|A| = b \Rightarrow A = b$ or $A = -b$
\item If absolute value = negative, no solution
\end{itemize}

\vspace{0.5cm}

\textbf{Rational Equations:}
\begin{itemize}
\item Find the LCD of all denominators
\item Multiply both sides by LCD
\item Solve the resulting equation
\item Always check solutions in the original equation!
\end{itemize}
\end{frame}

\end{document}
